\documentclass[10pt,a4paper]{amsart}
\usepackage[margin=19mm]{geometry}
\usepackage{amsmath,amssymb,mathtools}
\usepackage[scr=rsfs,cal=euler]{mathalfa}
\usepackage{xfrac}
\usepackage[unicode,hidelinks,bookmarks=false]{hyperref}
\newcommand{\ie}{i.e.\ }
\newcommand{\cf}{cf.\ }
\newcommand{\resp}{respectively\ }
\newcommand{\Subsection}{\subsection}
\providecommand{\nobreakdash}{\nobreak\hbox{-}}
\setlength{\emergencystretch}{2em}
\multlinegap0pt
\usepackage{amssymb}
\usepackage{mathtools}
\newtheorem{theorem}{Theorem}
\newtheorem{corollary}{Corollary}
\DeclareMathOperator{\Aut}{Aut}
\title{On base point freeness for rank one foliations}
\author{Paolo Cascini, Calum Spicer}
\date{Source: arXiv:2509.03109v2 (CC BY 4.0)}
\begin{document}
\maketitle

\noindent\emph{Source and licence.} On base point freeness for rank one foliations. Version arXiv:2509.03109v2 (CC BY 4.0), distributed by arXiv under the Creative Commons Attribution 4.0 International licence, http://creativecommons.org/licenses/by/4.0/, as recorded in the arXiv licence metadata for this exact version and shown on the abstract page https://arxiv.org/abs/2509.03109v2. Full licence notice: \texttt{licenses/arxiv-2509.03109-CC-BY-4.0.txt}.\par\smallskip

\noindent\textit{Editorial scope.} Counted unit: the article's corollary cor\_hyperbolicity (source0.tex lines 536--543). The article's complete original proof of it is delivered with it (source0.tex lines 544--551, one \texttt{proof} environment of 8 lines). No mathematics is written, repaired or re-derived: the statement and the proof are the article's own lines, in the article's own notation, and no retained line is substituted or re-flowed.  Retained uncounted as the source-local context the proof needs: lines 336--341.  Nothing was omitted from any retained span, so the package carries no omission record.  No retained line contains a citation, so the wrapper carries no bibliography and no bibliography entry.  No retained line contains a figure environment, an image-inclusion command or drawing code, so no image is delivered and the package declares no asset dependency.  Editorial formatting only, in addition: an AMS \texttt{amsart} preamble that restates the article's own declarations and macros used by the retained lines, namely the environments \texttt{theorem}, \texttt{corollary} on separate counters, together with the standard packages those lines need.  The retained environments print in this document as Theorem 1, Corollary 1 in order of appearance, whereas the article prints the same items with the numbering of its own full document.  The complete native source of the article as distributed in the arXiv e-print for this exact version is bundled unchanged as \texttt{sources/184391/source0.tex}.
\begin{theorem}\label{t_main}
Let $X$ be a normal projective $\mathbb Q$-factorial klt variety, let $\mathcal F$ be a rank one foliation on $X$ and let $\Delta \ge 0$ be a $\mathbb Q$-divisor such that $(\mathcal F, \Delta)$ is log canonical and 
$K_{\mathcal F}+\Delta \equiv 0$.  

Then $K_{\mathcal F}+\Delta \sim_{\mathbb Q}0$.
\end{theorem}

\begin{corollary}
\label{cor_hyperbolicity}
Let $X$ be a normal projective $\mathbb Q$-factorial klt variety, let $\mathcal F$ be a rank one foliation on $X$ such that $\mathcal F$ is log canonical and  
$K_{\mathcal F}\equiv 0$.

Then, for a general point $x \in X$ there exists a holomorphic map $f\colon \mathbb C \to X$ such that 
$x \in f(\mathbb C)$ and the image of $f$ is tangent to $\mathcal F$.
\end{corollary}

\begin{proof}
By Theorem \ref{t_main} $K_{\mathcal F}\sim_{\mathbb Q} 0$.  So, up to replacing $X$ by the index one cover associated to $K_{\mathcal F}$, we may assume that $K_{\mathcal F} \sim 0$.
Moreover, up to replacing $X$ by a functorial resolution of singularities, we may assume that $X$ is smooth.

Since $K_{\mathcal F} \sim 0$ we see that  $\mathcal F$ is generated by a global vector field $v \in H^0(X, T_X)$.  Since $H^0(X, T_X)$ is the Lie algebra of $\Aut (X)$, 
for a general point $x \in X$, we have the exponential map $\exp_x\colon H^0(X, T_X) \to X$ such
that $\exp_x(0) = x$.  We take $f\colon \mathbb C \to X$ to be the restriction of $\exp_x$ to the subspace $\mathbb Cv \subset H^0(X, T_X)$. 
\end{proof}

\end{document}
